\section{Success certificates, partial advances, and kill criteria}
\label{sec:decision-rules}

This section turns the preceding identities into a branchable
research plan.  The purpose is to make a failed route informative
without converting it into a claim about the whole parity problem.

\subsection{Route A: raw mass followed by fiber survival}

Route A succeeds at losses
\((\lambda_R,\lambda_{\rm fib},\lambda_S)\) only if all of the
following are proved for the literal fixed-\(h_0\) packet:
\begin{align}
 R&\ge X^{-\lambda_R+o(1)}N_0,
\label{eq:routeA-raw}\\
 M&\ge X^{-\lambda_{\rm fib}+o(1)}R,
\label{eq:routeA-fiber}\\
 S&\le X^{\lambda_S+o(1)}Q.
\label{eq:routeA-support}
\end{align}
Then
\[
 \lambda_D^{\rm cert}
 =2\lambda_R+2\lambda_{\rm fib}+\lambda_S.
\]
The verified support theorem supplies \(\lambda_S=0\).  If a
zero-mode transfer has gain \(\eta_Z\), require
\[
 \lambda_D^{\rm cert}\le2\eta_Z,
\]
and separately verify the strict complete loss budget
\eqref{eq:strict-endpoint}.  The certified energy loss is not
automatically added to \(\Lambda_{\rm phys}\); such an addition
requires the crosswalk described after
\eqref{eq:strict-endpoint}.

The fiber step may use the sharp magnitude certificate:
\[
 \sum_n(2r_{\max,n}-R_n)_+
 \ge X^{-\lambda_{\rm fib}+o(1)}R.
\]
If that fails, Route A remains open to a signed arithmetic lower
bound for \(M\).

\subsection{Route B: a positive pairing main term}

Choose a rigorously defined conjugation-closed pair family \(\cP\)
and write
\[
 C_{\eqc}=C_{\cP}+C_{\cR}.
\]
Route B succeeds if
\begin{align}
 C_{\cP}
 &\ge X^{-\lambda_P+o(1)}Q^3J^2,
\label{eq:routeB-main}\\
 |C_{\cR}|
 &=o(C_{\cP}),
\label{eq:routeB-rem}
\end{align}
or if a one-sided remainder estimate leaves a positive constant
fraction of \(C_{\cP}\).  Then
\[
 D=E_{\trip}+C_{\cP}+C_{\cR}
\ge X^{-\lambda_P+o(1)}Q^3J^2.
\]
The candidate family must preserve the exact determinant and all
physical orientations.  A pairing main term found only after
averaging over \(h\) or changing the packet is not a Route B
certificate.

\subsection{Publishable partial advances}

The following remain useful even when neither route reaches the final
endpoint:
\begin{enumerate}[leftmargin=2.2em]
 \item a rigorous positive raw exponent \(\lambda_R\);
 \item a nontrivial magnitude-survival exponent for a positive share
 of the raw mass;
 \item an exact multiplicity theorem for one pairing class;
 \item a positive \(C_{\cP}\) main term with a remainder too large for
 the endpoint;
 \item a localization theorem showing that a broad Chowla-hard
 remainder is confined to a small, explicitly weighted class;
 \item an exact negative census showing that a named symmetry or
 duplicate mechanism has insufficient mass.
\end{enumerate}
Each result must retain its \(\Lzero/\Lone/\Ltwo\) label and measured
loss.  A positive exponent larger than the available budget is
progress on the local gate, not completion of the physical route.

\subsection{Kill criteria}

A branch stops, or is relabeled purely exploratory, if any of the
following occurs:
\begin{enumerate}[leftmargin=2.2em]
 \item it proves only upper bounds for \(R\) or \(M\);
 \item the sharp magnitude envelope is zero on all but a negligible
 share of the raw mass, and the branch has no signed mechanism;
 \item a proposed pairing changes determinant sign, target,
 orientation, cutoff sign, or physical multiplicity;
 \item its gain appears only after averaging the fixed shift,
 determinant frequency, row slopes, masks, or incompatible packets;
 \item the required remainder theorem, uniformly stated, contains an
 unproved Chowla specialization and no geometry-specific
 localization;
 \item the certified energy loss violates
 \(\lambda_D^{\rm cert}\le 2\eta_Z\);
 \item the complete physical loss is at least \(1/400\), including
 equality.
\end{enumerate}

A killed branch supports only a branch-local negative conclusion:
for example, ``magnitude dominance does not certify fiber survival in
this packet.''  It does not prove a universal parity barrier or rule
out a different fixed-shift mechanism.

\subsection{Next theorem target}

The most informative next theorem is one of the following:
\begin{enumerate}[leftmargin=2.2em,label=(\roman*)]
 \item a raw nondegeneracy lower bound for \(R\);
 \item a census theorem showing
 \(\sum_nd_n\ge X^{-\lambda_{\rm fib}+o(1)}R\);
 \item an exact positive main term from duplicate or
 shared-target/shared-divisor pairings;
 \item a localization of the unpaired remainder to a family on which
 an existing mean-value theorem applies without introducing an
 unproved special case.
\end{enumerate}
A future executed census should decide which of (ii)--(iv) has enough
literal mass before a long analytic argument is attempted.  The
schema in \cref{sec:census-protocol} does not itself make that
decision.
